\documentclass[11pt]{article}
\usepackage[margin=1in]{geometry}
\usepackage{amsmath,amssymb,amsthm,mathtools}
\usepackage{graphicx}
\usepackage{microtype}
\usepackage[hidelinks]{hyperref}

\newtheorem{theorem}{Theorem}
\newtheorem{lemma}{Lemma}
\newtheorem{remark}{Remark}
\newcommand{\Conv}{\operatorname{Conv}^{\infty}}
\newcommand{\VN}{\operatorname{VN}}
\newcommand{\F}{\mathcal F}
\newcommand{\one}{\mathbf 1}

\title{Bounded Convolution Kernels on Locally Compact Abelian Groups\\
\large A complete LCA subcase of questions of Forrest--Spronk--Wood}
\author{Autonomous proof candidate for human review}
\date{11 August 2026}

\begin{document}
\maketitle

\begin{center}
\fbox{\parbox{0.91\textwidth}{\textbf{Status.}
Candidate partial result, likely valid.  The theorem completely settles two
of the source questions for locally compact abelian groups.  It does not settle
the questions for arbitrary locally compact groups and does not fully describe
the closure in the negative cases.}}
\end{center}

\section{Source question}

Let $G$ be a locally compact group.  Forrest, Spronk, and Wood
\cite{FSW} write $\Conv(G)$ for the bounded measurable functions $\varphi$
whose left-convolution operator extends boundedly to $L^2(G)$, and denote that
operator by $\Lambda(\varphi)$.  On page 8, at the end of Section 2.1, they
ask for which groups
\[
  \Conv(G)\subseteq L^2(G),
\]
when $\overline{\Lambda(\Conv(G))}^{\|\cdot\|}$ is a $*$-subalgebra of
$\VN(G)$, and how large that closure is.

\begin{figure}[ht]
  \centering
  \includegraphics[width=\textwidth]{figures/open_problem_crop.png}
  \caption{The questions in the original paper, source PDF page 8.}
\end{figure}

The source observes that compact and discrete groups have both requested
properties.  The result below proves that, within the entire LCA class, there
are no other positive cases.

\section{Proof intuition}

Under the Fourier transform on an abelian group, convolution becomes
multiplication.  Thus a bounded function is a bounded $L^2$-convolver exactly
when its distributional Fourier transform is a bounded function.  The problem
becomes an uncertainty question: when can both a function and its Fourier
transform be bounded without the original function lying in $L^2$?

The LCA structure theorem gives an open subgroup $\mathbb R^d\times K$ with
$K$ compact.  If $d>0$, a quadratic chirp has a bounded, unimodular Fourier
transform and infinite $L^2$ norm.  Multiplying its convolution operator by
its adjoint produces the singular projection obtained by averaging over $K$;
the distance theorem in the source paper keeps that projection one unit away
from every bounded-function convolver.

If $d=0$, then $K$ is compact open.  Put different characters of $K$ on
different cosets of $K$.  Orthogonality makes the Fourier transform bounded on
disjoint fibers, although the kernel has infinite $L^2$ norm.  The product
with its adjoint is the indicator of the union of those fibers.  Any uniform
approximation to this projection would, after averaging along a dual fiber,
give Fourier coefficients of an $L^1(K)$ function that remain close to $1$
along infinitely many distinct characters.  Riemann--Lebesgue forbids this.

\section{Fourier formulation}

Assume henceforth that $G$ is LCA.  Choose dual Haar measures so that
$\F:L^2(G)\to L^2(\widehat G)$ is unitary.

\begin{lemma}[Multiplier criterion]\label{lem:multiplier}
For $\varphi\in L^\infty(G)$, the following are equivalent:
\begin{enumerate}
\item $\varphi\in\Conv(G)$;
\item the distributional Fourier transform $\widehat\varphi$ is represented
by some $m_\varphi\in L^\infty(\widehat G)$.
\end{enumerate}
In this case
\[
 \F\Lambda(\varphi)\F^{-1}=M_{m_\varphi},\qquad
 \|\Lambda(\varphi)\|=\|m_\varphi\|_\infty.
\]
Moreover, if $U$ is an open subgroup, extension by zero of a kernel on $U$
corresponds to pulling its multiplier back by the restriction map
$r:\widehat G\to\widehat U$.
\end{lemma}

\begin{proof}
A bounded translation-invariant operator on $L^2(G)$ is diagonalized by the
Plancherel transform and hence is multiplication by an $L^\infty$ function.
Testing first on compactly supported Fourier test functions identifies this
multiplier with the distributional Fourier transform of the convolution
kernel.  The converse follows from Plancherel and the distributional
convolution identity.  For an open subgroup, integration of the zero extension
is integration over $U$, so the same test-function calculation gives pullback
along $r$.
\end{proof}

\section{Classification theorem}

\begin{theorem}\label{thm:main}
Let $G$ be a locally compact abelian group.
\begin{enumerate}
\item $\Conv(G)\subseteq L^2(G)$ if and only if $G$ is compact or discrete.
\item The operator-norm closure
\[
 \mathcal C_G:=\overline{\Lambda(\Conv(G))}^{\|\cdot\|}\subseteq\VN(G)
\]
is a $*$-subalgebra if and only if $G$ is compact or discrete.
\end{enumerate}
If $G$ is neither compact nor discrete, the proof constructs
$\varphi\in\Conv(G)\setminus L^2(G)$ such that
$\Lambda(\varphi)\Lambda(\varphi)^*$ has distance exactly $1$ from
$\Lambda(\Conv(G))$ after normalizing $\|\Lambda(\varphi)\|=1$.
\end{theorem}

\begin{proof}
\textbf{The positive cases.}
If $G$ is compact, Haar measure is finite, so $L^\infty(G)\subseteq L^2(G)$.
Also $\Conv(G)=L^\infty(G)$, and bounded functions are dense in $L^1(G)$;
hence $\mathcal C_G=C_r^*(G)$.

If $G$ is discrete and $\varphi\in\Conv(G)$, then
$\varphi=\Lambda(\varphi)\delta_e\in\ell^2(G)$.  Conversely every element of
$\VN(G)$ is convolution by the $\ell^2$ vector obtained by applying it to
$\delta_e$; this vector is bounded.  Therefore
$\Lambda(\Conv(G))=\VN(G)$.  Both assertions follow in these two cases.

\medskip
\textbf{Structural reduction.}
The LCA structure theorem gives an open subgroup
\[
 U\cong\mathbb R^d\times K
\]
with $K$ compact \cite{HR}.  Suppose that $G$ is neither compact nor discrete.
We split according to $d$.

\medskip
\textbf{Case 1: $d>0$.}
Normalize Haar measure on $K$ to have mass $1$.  On $U$ set
\[
 \varphi_U(x,k)=e^{\pi i|x|^2},
\]
and extend $\varphi_U$ by zero to $G$.  The Fresnel identity, in tempered
distributions, gives
\[
 \widehat{e^{\pi i|\cdot|^2}}(\xi)
   =\zeta_d e^{-\pi i|\xi|^2},\qquad |\zeta_d|=1.
\]
Consequently, on
$\widehat U\cong\mathbb R^d\times\widehat K$,
\[
 m_U(\xi,\gamma)
 =\zeta_d e^{-\pi i|\xi|^2}\one_{\{1_{\widehat K}\}}(\gamma).
\]
Lemma~\ref{lem:multiplier} shows that the zero extension $\varphi$ belongs to
$\Conv(G)$ with convolution norm $1$.  Since $|\varphi_U|=1$ on a set of
infinite Haar measure, $\varphi\notin L^2(G)$.

Let $P=\Lambda(\varphi)\Lambda(\varphi)^*$.  Its multiplier is the pullback
of $|m_U|^2$, so
\[
 P=\int_K\lambda(k)\,dk,
\]
the orthogonal projection obtained by averaging over $K$.  In particular
$P\in\VN_K(G)$ and $\|P\|=1$.  The subgroup $K$ is closed and non-open,
because $U/K\cong\mathbb R^d$.  Every LCA group is amenable, so $K$ admits a
contractive approximate indicator in $G$.  The distance theorem
\cite[Theorem~3.1]{FSW} therefore yields
\[
 \operatorname{dist}\bigl(P,\Lambda(\Conv(G))\bigr)=\|P\|=1.
\]
Thus $P\notin\mathcal C_G$, although both factors defining $P$ lie in
$\Lambda(\Conv(G))$.  Hence $\mathcal C_G$ is not an algebra.

\medskip
\textbf{Case 2: $d=0$.}
Now $U=K$ is compact open.  Since $G$ is nondiscrete, $K$ is infinite; since
$G$ is noncompact, the discrete quotient $G/K$ is infinite.  Its annihilator
\[
 K^\perp=\{\chi\in\widehat G:\chi|_K=1\}
\]
is compact, and the restriction map
\[
 r:\widehat G\longrightarrow\widehat K
\]
is onto, with fibers the cosets of $K^\perp$.

Choose distinct cosets $x_n+K$ in $G/K$ and distinct characters
$\gamma_n\in\widehat K$, for $n\geq1$.  Define a bounded measurable function
by
\[
 \varphi(x_n+k)=\gamma_n(k)\quad(k\in K),
 \qquad
 \varphi=0\quad\text{off }\bigcup_{n\geq1}(x_n+K).
\]
Each occupied coset contributes $1$ to the squared $L^2$ norm, and hence
$\varphi\notin L^2(G)$.

For $N\geq1$, let $\varphi_N$ be the truncation to the first $N$ occupied
cosets.  Then $\varphi_N\in L^1(G)\cap L^\infty(G)$, and character
orthogonality on $K$ gives
\[
 \widehat{\varphi_N}(\chi)=
 \begin{cases}
   \overline{\chi(x_n)},&r(\chi)=\gamma_n\text{ for some }n\leq N,\\
   0,&r(\chi)\notin\{\gamma_n:1\leq n\leq N\}.
 \end{cases}
\]
The displayed multipliers have pairwise disjoint fiber supports and converge
pointwise to
\[
 m(\chi)=
 \begin{cases}
   \overline{\chi(x_n)},&r(\chi)=\gamma_n\text{ for some }n,\\
   0,&r(\chi)\notin\{\gamma_n:n\geq1\}.
 \end{cases}\label{eq:multiplier}
\]
In particular $m\in L^\infty(\widehat G)$ and $\|m\|_\infty=1$.  If
$f\in L^1(G)\cap L^2(G)$, dominated convergence gives
$\varphi_N*f\to\varphi*f$ pointwise, while Plancherel and dominated
convergence give
\[
 \varphi_N*f\longrightarrow\F^{-1}(m\widehat f)
 \quad\text{in }L^2(G).
\]
After passage to an almost-everywhere convergent subsequence, the two limits
agree.  Hence convolution by $\varphi$ extends to the multiplier $M_m$, and
$\varphi\in\Conv(G)$.

It remains to show that the closure is not an algebra.  Set
\[
 E=r^{-1}\bigl(\{\gamma_n:n\geq1\}\bigr).
\]
Then the multiplier of
$P:=\Lambda(\varphi)\Lambda(\varphi)^*$ is $|m|^2=\one_E$.
We claim
\[
 \operatorname{dist}\bigl(P,\Lambda(\Conv(G))\bigr)=1.\label{eq:distance}
\]
The upper bound follows by using the zero operator.  For the lower bound, let
$\theta\in\Conv(G)$ have multiplier $a\in L^\infty(\widehat G)$, and suppose
for contradiction that
\[
 \|a-\one_E\|_\infty=\delta<1.
\]
Let $\mu$ be Haar probability measure on the compact subgroup $K^\perp$ and
average along its cosets:
\[
 A=a*\mu.
\]
Convolution by $\mu$ is contractive on $L^\infty$, and $E$ is a union of
$K^\perp$-cosets, so
\[
 \|A-\one_E\|_\infty\leq\delta.\label{eq:avg-close}
\]
The inverse Fourier transform of $\mu$ is $\one_K$.  Therefore $A$ is the
Fourier transform of $\theta\one_K$.  Since
$\theta\one_K\in L^1(K)$, $A$ factors through $\widehat K$ and its values are
the Fourier coefficients
\[
 A(\widetilde\gamma)
 =\int_K\theta(k)\overline{\gamma(k)}\,dk,
 \qquad r(\widetilde\gamma)=\gamma.
\]
By Riemann--Lebesgue these coefficients form an element of
$c_0(\widehat K)$.  In particular their values along the distinct
$\gamma_n$ tend to zero.  But \eqref{eq:avg-close} says that their absolute
values are at least $1-\delta>0$ for every $n$, a contradiction.  This proves
\eqref{eq:distance}, and again $P$ is the product of two members of
$\Lambda(\Conv(G))$ but lies outside $\mathcal C_G$.

The two structural cases exhaust every LCA group that is neither compact nor
discrete, proving both equivalences.
\end{proof}

\begin{remark}[Concrete compact-open example]
For $G=\mathbb Z\times\mathbb T$, the compact-open construction can be seen
without any structure theory: $\varphi(n,t)=e^{2\pi i nt}$ is bounded and not
in $L^2(G)$, while its Fourier transform on
$\mathbb T\times\mathbb Z$ is again unimodular after exchanging the two
variables.
\end{remark}

\section{Verification and limitations}

The proof was checked branch by branch against the source definitions.
The main verifier targets were: (i) the finite-truncation and dominated-limit
argument leading to \eqref{eq:multiplier}; (ii) the identity turning dual-fiber
averaging into multiplication by $\one_K$; and (iii) the hypotheses of the
source distance theorem in Case 1.  All three checks pass.  No computation is
used as proof.

This result is complete only in the LCA subclass.  It does not classify
arbitrary locally compact groups, and it does not give a full intrinsic
description of $\mathcal C_G$ in the negative cases.  A bounded local-corpus
and web/arXiv search found the source paper and general literature on
convolvers and pseudomeasures, but no paper stating these two LCA
equivalences.  Because the ingredients are classical and older literature may
use different terminology, novelty confidence is moderate rather than high.

\section{Human review recommendation}

Send to a harmonic analyst for review.  The proof is short enough to audit
directly.  The compact-open construction and the fiber-average
Riemann--Lebesgue obstruction deserve the closest check; bibliographic review
should search older pseudomeasure terminology before a public novelty claim.

\begin{thebibliography}{9}
\bibitem{FSW}
B. E. Forrest, N. Spronk, and P. J. Wood,
\emph{Operator Segal algebras in Fourier algebras},
Studia Math. \textbf{179} (2007), no.~3, 277--295;
arXiv:math/0601509.

\bibitem{HR}
E. Hewitt and K. A. Ross,
\emph{Abstract Harmonic Analysis, Vol. I},
2nd ed., Springer, 1979.
\end{thebibliography}

\end{document}
